\begin{lemma}[Smooth approximation of an $L^1(\gamma)$ datum, and the three-term splitting of every truncated Riesz integral]
Fix $n \geq 1$, let $\lambda$ be Lebesgue measure on $\mathbb{R}^n$, let
$\gamma_n$ be the Gaussian measure of \noderef{GaussMeasure}, and let
$K_{\gamma,j}$ be the kernel of \noderef{GaussRieszKernel}.  For a
$\gamma_n$-integrable $h : \mathbb{R}^n \to \mathbb{R}$, a coordinate
$j \in \{1,\ldots,n\}$, a point $x \in \mathbb{R}^n$ and a real
$\varepsilon > 0$ write
\[
  T_{j,\varepsilon}h(x)
    = \int_{|x-y|>\varepsilon}K_{\gamma,j}(x,y)\,h(y)\,dy ,
\]
the $y$-integral against $\lambda$; for an element of $L^1(\gamma_n)$ this is
read through its canonical representative.

Let $f \in L^1(\gamma_n)$ and let $k \in \mathbb{N}$.  Then there exist a
function $\varphi : \mathbb{R}^n \to \mathbb{R}$ and an element
$g \in L^1(\gamma_n)$ such that
\begin{enumerate}
\item[(i)] $\varphi$ is smooth, has compact support, and is
$\gamma_n$-integrable;
\item[(ii)] $\displaystyle\int_{\mathbb{R}^n}|g|\,d\gamma_n
  \;<\; \frac1{k+1}$, the integral of the canonical representative of $g$; and
\item[(iii)] $f = \varphi + g$ holds $\gamma_n$-almost everywhere, as an
identity between the canonical representative of $f$, the function $\varphi$,
and the canonical representative of $g$; equivalently $f = [\varphi] + g$ in
$L^1(\gamma_n)$, where $[\varphi]$ is the class of $\varphi$, which exists by
(i); and
\item[(iv)] for every coordinate $j$, every $x \in \mathbb{R}^n$ and every
real $\varepsilon > 0$, each of the three functions
\[
  y \mapsto K_{\gamma,j}(x,y)f(y), \qquad
  y \mapsto K_{\gamma,j}(x,y)\varphi(y), \qquad
  y \mapsto K_{\gamma,j}(x,y)g(y)
\]
is Lebesgue integrable on $\{y : |x-y| > \varepsilon\}$, and
\[
  T_{j,\varepsilon}f(x)
    \;=\; T_{j,\varepsilon}\varphi(x) + T_{j,\varepsilon}g(x) .
\]
\end{enumerate}

\emph{Why $g$ is an element of $L^1(\gamma_n)$ and not the function
$f - \varphi$.}  The intended consumer feeds $g$ to a weak-type estimate whose
input is an $L^1(\gamma_n)$ element, so what it needs is a class, not a
representative; and it needs the splitting (iv) stated for \emph{that class's
own} canonical representative, since that is the function whose truncations
and maximal truncation the estimate controls.  Stating (iv) with $g$ rather
than with $f - \varphi$ is therefore not a convenience: it discharges here,
once, the reconciliation between the class $f - [\varphi]$ and the pointwise
difference $f - \varphi$, which otherwise has to be carried through every later
step.  The two agree $\gamma_n$-almost everywhere and hence, by clause (4) of
\noderef{GaussLebesgueBallComparison}, $\lambda$-almost everywhere, which is
what makes the substitution legitimate inside a Lebesgue integral.

\emph{Why (iii) is stated between representatives rather than between
classes.}  The two formulations are equivalent --- two elements of
$L^1(\gamma_n)$ whose representatives agree $\gamma_n$-almost everywhere are
equal --- and the representative form is the one a consumer can use without
first naming the class $[\varphi]$, which needs the membership (i) as a datum.
A consumer that wants the class form reads it off: (iii) says that $g$ and
$f - \varphi$ agree $\gamma_n$-almost everywhere, so $g = f - [\varphi]$ in
$L^1(\gamma_n)$, and hence
$\|g\|_{L^1(\gamma_n)} = \|f - [\varphi]\|_{L^1(\gamma_n)}$ and, for any linear
$T$ on $L^1(\gamma_n)$, $Tg = Tf - T[\varphi]$.

\emph{The finiteness in (iv) is at every point, with no exceptional set.}
Each of the three integrands is asserted integrable on
$\{|x-y| > \varepsilon\}$ for \emph{every} $x$ and every $\varepsilon > 0$, so
each truncation is a finite real number everywhere, and the identity in (iv)
is an identity between finite reals rather than an almost-everywhere statement.
This is available because the far majorant of
\noderef{GaussRieszKernelFarMajorant} is a genuine pointwise kernel bound off
the diagonal; no maximal-function or almost-everywhere finiteness input is used
or needed.

\emph{What the bound in (ii) does and does not say.}  The constant $1/(k+1)$ is
the only role $k$ plays: the statement is an approximation statement at the
single accuracy $1/(k+1)$, and $\varphi$ and $g$ both depend on $k$.  Nothing
is claimed about how $\varphi$ or $g$ vary with $k$, and in particular no
uniformity in $k$ and no convergence as $k \to \infty$ is asserted; a consumer
that wants a sequence applies the lemma once for each $k$.  Note also that
$\varphi$ and $g$ do not depend on the coordinate $j$: $j$ is quantified
inside (iv), after both have been produced.
\end{lemma}
\begin{proof}
\emph{Step 1: the approximant and the class $g$.}
Apply \noderef{SmoothCompactDenseL1Gauss} to $f$ with
$\delta = 1/(k+1) > 0$.  It produces a function
$\varphi : \mathbb{R}^n \to \mathbb{R}$ which is smooth, has compact support,
is $\gamma_n$-integrable, and satisfies
\begin{equation*}
  \bigl\|\,f - \varphi\,\bigr\|_{L^1(\gamma_n)} < \frac1{k+1} , \tag{1}
\end{equation*}
where $f - \varphi$ means the pointwise difference of the canonical
representative of $f$ and $\varphi$, and the norm is the $L^1(\gamma_n)$ norm
of that function.  This gives (i).

Let $[\varphi] \in L^1(\gamma_n)$ be the class of $\varphi$, which exists
because $\varphi$ is $\gamma_n$-integrable, and put
\[
  g = f - [\varphi] \;\in\; L^1(\gamma_n) .
\]
The canonical representative of $g$ agrees $\gamma_n$-almost everywhere with
the pointwise difference $f - \varphi$: the representative of a difference of
classes agrees almost everywhere with the difference of the representatives,
and the representative of $[\varphi]$ agrees almost everywhere with $\varphi$.
Hence
\begin{equation*}
  g = f - \varphi \quad \gamma_n\text{-almost everywhere.} \tag{2}
\end{equation*}
Rearranging $(2)$ pointwise at each of the $\gamma_n$-almost every $y$
at which it holds gives $f(y) = \varphi(y) + g(y)$ there, which is (iii).

Since the $L^1(\gamma_n)$ norm of a function is unchanged by a
$\gamma_n$-almost-everywhere modification, $(2)$ and $(1)$ give
$\|g\|_{L^1(\gamma_n)} < 1/(k+1)$; and
$\|g\|_{L^1(\gamma_n)} = \int|g|\,d\gamma_n$, the integral of the canonical
representative.  That is (ii).

\emph{Step 2: every truncation is finite, and the splitting.}
Fix $j$, $x \in \mathbb{R}^n$ and a real $\varepsilon > 0$, and put
$R_x = \max(|x|,1)$, so that $R_x > 0$ and $|x| \leq R_x$.  Apply clause (1) of
\noderef{GaussRieszKernelFarMajorant} with $R := R_x > 0$ and
$\rho := \varepsilon > 0$.  Its second part states that for every
$h \in L^1(\gamma_n)$ --- more precisely for every $h$ that is
$\gamma_n$-integrable --- and every $x'$ with $|x'| \leq R_x$, the function
$y \mapsto K_{\gamma,j}(x',y)h(y)$ is Lebesgue integrable on
$\{y : |x'-y| > \varepsilon\}$.  Taking $x' = x$, which is legitimate because
$|x| \leq R_x$, and taking $h$ in turn to be the canonical representative of
$f$, then $\varphi$, then the canonical representative of $g$ --- all three
being $\gamma_n$-integrable, the first and third as representatives of elements
of $L^1(\gamma_n)$ and the second by (i) --- gives the three integrability
assertions of (iv).  In particular each of
$T_{j,\varepsilon}f(x)$, $T_{j,\varepsilon}\varphi(x)$ and
$T_{j,\varepsilon}g(x)$ is a finite real number.

For the identity, note first that by clause (4) of
\noderef{GaussLebesgueBallComparison} a subset of $\mathbb{R}^n$ is
$\gamma_n$-null if and only if it is $\lambda$-null, so the
$\gamma_n$-almost-everywhere identity $(2)$ is also a
$\lambda$-almost-everywhere identity:
\begin{equation*}
  g(y) = f(y) - \varphi(y) \quad \text{for } \lambda\text{-almost every } y .
  \tag{3}
\end{equation*}
Consequently, for $\lambda$-almost every $y$,
\[
  K_{\gamma,j}(x,y)f(y)
    = K_{\gamma,j}(x,y)\varphi(y) + K_{\gamma,j}(x,y)\bigl(f(y)-\varphi(y)\bigr)
    = K_{\gamma,j}(x,y)\varphi(y) + K_{\gamma,j}(x,y)g(y) ,
\]
the first equality being the distributive law at the point $y$ and the second
substituting $(3)$.  The set $\{y : |x-y| > \varepsilon\}$ is Lebesgue
measurable, being the strict superlevel set of the continuous function
$y \mapsto |x-y|$, so two functions agreeing $\lambda$-almost everywhere have
the same integral over it; hence
\[
  T_{j,\varepsilon}f(x)
    = \int_{|x-y|>\varepsilon}
        \bigl(K_{\gamma,j}(x,y)\varphi(y) + K_{\gamma,j}(x,y)g(y)\bigr)dy .
\]
Both summands are Lebesgue integrable on $\{y : |x-y| > \varepsilon\}$ by the
previous paragraph, so the integral of the sum is the sum of the integrals:
\[
  T_{j,\varepsilon}f(x)
    = \int_{|x-y|>\varepsilon}K_{\gamma,j}(x,y)\varphi(y)\,dy
      + \int_{|x-y|>\varepsilon}K_{\gamma,j}(x,y)g(y)\,dy
    = T_{j,\varepsilon}\varphi(x) + T_{j,\varepsilon}g(x) .
\]
As $j$, $x$ and $\varepsilon > 0$ were arbitrary, this is (iv).
\end{proof}
